\documentclass[11pt,a4paper]{article}
\usepackage[T1]{fontenc}
\usepackage{lmodern}
\usepackage[utf8]{inputenc}
\usepackage[polish]{babel}
\usepackage{amsmath,amssymb,hyperref}
\usepackage[margin=2.25cm]{geometry}
\setlength{\parindent}{0pt}
\setlength{\parskip}{4pt}
\newcommand{\topic}[3]{%
  \noindent\begin{minipage}[t][.42\textheight][t]{\linewidth}
  {\large\bfseries #1. #2}\par\smallskip
  \small #3
  \end{minipage}\par\vspace{.35cm}
}
\title{Zagadnienia na egzamin licencjacki z fizyki}
\author{Opracowanie dydaktyczne}
\date{\today}
\begin{document}
\maketitle
Podstawą zakresu jest
\href{https://wfa.uwr.edu.pl/wp-content/uploads/sites/216/2025/10/Pytania-lic-Fizyka-I-st.pdf}
{oficjalna lista pytań Wydziału Fizyki i Astronomii UWr}.
Każdemu zagadnieniu przydzielono około pół strony układu.
Tekst jest skróconym konspektem; wymaga rozwinięcia przykładów
i wyprowadzeń przed wykorzystaniem jako pełnego opracowania
egzaminacyjnego.
\newpage
\section*{I. Fizyka ogólna}
\topic{1}{Zasady dynamiki Newtona, pęd i jego zachowanie}{
W układzie inercjalnym ciało zachowuje stan spoczynku lub ruchu
jednostajnego prostoliniowego, jeżeli siła wypadkowa jest zerowa.
Druga zasada ma ogólną postać $\mathbf F_{\rm wyp}=d\mathbf p/dt$;
przy stałej masie redukuje się do $m\mathbf a$.
Trzecia zasada mówi, że siły wzajemnego oddziaływania dwóch ciał
mają przeciwne zwroty i tę samą wartość, ale działają na różne ciała.

Dla układu ciał sumujemy pędy: $\mathbf P=\sum_i\mathbf p_i$.
Siły wewnętrzne znoszą się parami w bilansie całego układu,
więc $d\mathbf P/dt=\mathbf F_{\rm zewn}$. Gdy zewnętrzny impuls
jest pomijalny, $\mathbf P$ pozostaje stały. Dotyczy to także
zderzeń niesprężystych, w których energia kinetyczna może maleć.
Przed zastosowaniem zasady trzeba jasno określić granice układu
i sprawdzić, czy impuls sił zewnętrznych jest rzeczywiście mały.}
\topic{2}{Układy inercjalne i nieinercjalne, siły bezwładności}{
Układ inercjalny to taki, w którym swobodne ciało porusza się
jednostajnie po prostej. Układ poruszający się względem niego
ze stałą prędkością także jest inercjalny. W układzie
przyspieszającym równanie Newtona wymaga dodania sił pozornych.
Dla układu przesuwającego się z przyspieszeniem $\mathbf A$
pojawia się siła $-m\mathbf A$.

W układzie obracającym się z prędkością kątową $\boldsymbol\Omega$
występują między innymi siła Coriolisa
$-2m\boldsymbol\Omega\times\mathbf v'$ i odśrodkowa
$-m\boldsymbol\Omega\times(\boldsymbol\Omega\times\mathbf r)$.
Jeżeli $\boldsymbol\Omega$ się zmienia, dochodzi siła Eulera.
Siły bezwładności nie wynikają z oddziaływania z innym ciałem;
są skutkiem wyboru współrzędnych. Przykładem jest odchylenie
toru długiego ruchu na obracającej się Ziemi.}
\newpage
\topic{3}{Praca, energia i siły zachowawcze}{
Praca siły na drodze $C$ jest całką
$W=\int_C\mathbf F\cdot d\mathbf r$.
Z twierdzenia o pracy i energii kinetycznej wynika
$W_{\rm wyp}=\Delta K$, gdzie $K=mv^2/2$ dla ruchu
nierelatywistycznego. Siła zachowawcza ma pracę niezależną
od drogi i można ją zapisać jako $\mathbf F=-\nabla U$.
Przykładami są grawitacja Newtonowska i idealna sprężyna.

Gdy działają tylko siły zachowawcze, $E=K+U$ jest stałe.
Jeśli działa tarcie, $\Delta(K+U)=W_{\rm niezach}$; utracona
energia mechaniczna przechodzi zwykle w energię wewnętrzną.
W zadaniu trzeba określić poziom zerowy energii potencjalnej
i rozróżnić zachowanie energii całkowitej od zachowania
energii mechanicznej. Moc chwilowa wynosi
$P=dW/dt=\mathbf F\cdot\mathbf v$.}
\topic{4}{Dynamika bryły sztywnej i moment pędu}{
Bryła sztywna zachowuje odległości między punktami.
Jej ruch można rozłożyć na translację środka masy oraz obrót.
Dla obrotu wokół ustalonej osi moment bezwładności jest
$I=\sum_i m_i r_{\perp i}^2$, a energia kinetyczna obrotu
$K_{\rm obr}=I\omega^2/2$. Moment siły to
$\boldsymbol\tau=\mathbf r\times\mathbf F$.

Równanie dynamiki obrotowej ma postać
$d\mathbf L/dt=\boldsymbol\tau_{\rm zewn}$.
Przy zerowym zewnętrznym momencie siły moment pędu
$\mathbf L$ jest zachowany. Dla obrotu wokół osi głównej
$L=I\omega$, więc zmniejszenie $I$ zwiększa $\omega$.
W ogólnym ruchu przestrzennym $\mathbf L$ i
$\boldsymbol\omega$ nie muszą być równoległe, ponieważ
moment bezwładności jest tensorem.}
\newpage
\topic{5}{Oscylator harmoniczny, tłumienie, wymuszenie i rezonans}{
Idealny oscylator spełnia $m\ddot x+kx=0$.
Rozwiązanie ma postać $x=A\cos(\omega_0t+\phi)$,
gdzie $\omega_0=\sqrt{k/m}$. Po dodaniu oporu lepkiego
$b\dot x$ otrzymujemy $m\ddot x+b\dot x+kx=0$.
W zależności od $b^2-4mk$ odpowiedź jest oscylacyjna
tłumiona, krytyczna albo aperiodyczna.

Przy wymuszeniu $F_0\cos\Omega t$ rozwiązanie zawiera
zanikający stan przejściowy i stan ustalony. Jego amplituda to
$A(\Omega)=F_0/\sqrt{(k-m\Omega^2)^2+b^2\Omega^2}$.
Dla słabego tłumienia maksimum występuje blisko
$\omega_0$; dokładne położenie maksimum wychylenia zależy
od $b$. W rezonansie energia dostarczana w kolejnych okresach
kompensuje straty. Należy odróżniać rezonans amplitudy
wychylenia od rezonansu mocy.}
\topic{6}{Grawitacja Newtona, prawa Keplera i energia potencjalna}{
Dwie masy punktowe przyciągają się siłą
$F=Gm_1m_2/r^2$ skierowaną wzdłuż prostej łączącej ciała.
Przy zerze energii w nieskończoności
$U(r)=-Gm_1m_2/r$. Ujemna energia całkowita w problemie
dwóch ciał odpowiada orbicie związanej.

Pierwsze prawo Keplera mówi o elipsie ze Słońcem
w jednym ognisku. Drugie, o równych polach zakreślanych
w równych czasach, wynika z zachowania momentu pędu.
Trzecie daje $T^2=4\pi^2a^3/[G(M+m)]$ dla półosi wielkiej
$a$ w problemie dwóch ciał. Prawo odwrotnych kwadratów
wyjaśnia te zależności przy założeniu izolowanego układu.
W pobliżu powierzchni Ziemi przy małych zmianach wysokości
można stosować przybliżenie $U\simeq mgh$.}
\newpage
\topic{7}{Statyka i dynamika płynów}{
W cieczy spoczywającej ciśnienie rośnie z głębokością:
$p=p_0+\rho gh$. Prawo Pascala opisuje przekazywanie
zmiany ciśnienia w cieczy zamkniętej. Prawo Archimedesa
daje siłę wyporu równą ciężarowi wypartego płynu,
$F_w=\rho_{\mathrm{plyn}}gV_{\rm wyparte}$.

Dla ustalonego przepływu nieściśliwego równanie ciągłości
ma postać $Av=\mathrm{const}$ w pojedynczym przewodzie.
Wzdłuż linii prądu płynu idealnego zachodzi
$p+\rho v^2/2+\rho gz=\mathrm{const}$.
Równanie Bernoulliego wynika z bilansu energii, lecz wymaga
braku istotnych strat lepkości i braku pomp między
porównywanymi punktami. W praktyce przepływ turbulentny
i tarcie wymagają dodatkowego członu strat.}
\topic{8}{Transformacje Galileusza i Lorentza}{
Transformacja Galileusza dla względnej prędkości $V$
ma postać $x'=x-Vt$, $t'=t$. Zakłada wspólny czas i działa
dobrze, gdy prędkości są małe względem $c$.
Transformacja Lorentza to
$x'=\gamma(x-Vt)$,
$t'=\gamma(t-Vx/c^2)$, gdzie
$\gamma=(1-V^2/c^2)^{-1/2}$.

Z niej wynikają dylatacja czasu własnego
$\Delta t=\gamma\Delta\tau$ i skrócenie długości
w kierunku ruchu $L=L_0/\gamma$. Prędkości wzdłuż
tej osi składają się według
$u'=(u-V)/(1-uV/c^2)$. Pęd relatywistyczny
to $\mathbf p=\gamma m\mathbf v$, energia całkowita
$E=\gamma mc^2$, a kinetyczna $K=(\gamma-1)mc^2$.
W granicy $v\ll c$ wracamy do $mv^2/2$.}
\newpage
\topic{9}{Gaz doskonały i jego przemiany}{
Model gazu doskonałego zakłada cząstki o zaniedbywalnej
objętości i oddziaływaniach poza zderzeniami.
Równanie stanu $pV=nRT=Nk_BT$ wiąże wielkości
makroskopowe. Temperatura jest związana ze średnią
energią ruchu cząstek, a ciśnienie wynika z przekazu
pędu ścianom.

W przemianie izotermicznej gazu doskonałego $pV$
jest stałe. W izochorycznej stała jest objętość,
w izobarycznej ciśnienie. Dla odwracalnej przemiany
adiabatycznej $pV^\gamma=\mathrm{const}$, gdzie
$\gamma=C_p/C_V$; wzór nie opisuje dowolnego procesu
z izolacją cieplną, jeśli jest on silnie nieodwracalny.
Praca wykonana przez gaz wynosi
$W=\int p\,dV$, a pole pod wykresem $p(V)$
ma interpretację energetyczną.}
\topic{10}{Pierwsza zasada termodynamiki i ekwipartycja}{
Pierwsza zasada to bilans energii wewnętrznej.
Przy konwencji, że $W$ jest pracą wykonaną przez układ,
$\Delta U=Q-W$. Dla gazu doskonałego energia wewnętrzna
zależy tylko od temperatury; w przemianie izotermicznej
$\Delta U=0$. W przemianie adiabatycznej $Q=0$,
więc rozprężanie kosztem pracy obniża $U$.

Klasyczna zasada ekwipartycji przyznaje średnio
$k_BT/2$ każdemu niezależnemu kwadratowemu składnikowi
energii w równowadze. Dla gazu jednoatomowego
$U=3Nk_BT/2$ i $C_V=3Nk_B/2$.
Nie wolno jednak bezkrytycznie stosować jej do wszystkich
stopni swobody przy niskich temperaturach: wzbudzenia
rotacyjne i drganiowe mogą być zamrożone kwantowo.
Rozróżniaj pojemność cieplną całego układu od molowej.}
\newpage
\topic{11}{Druga zasada, ciepło zredukowane i entropia}{
Druga zasada wyznacza kierunek przemian: w izolowanym
układzie całkowita entropia nie maleje.
Dla odwracalnej drogi $\Delta S=\int\delta Q_{\rm rev}/T$;
iloraz $\delta Q/T$ nazywa się ciepłem zredukowanym.
Entropia $S$ jest funkcją stanu, choć samo ciepło
zależy od drogi.

Nierówność Clausiusa
$\oint\delta Q/T\leq0$ odróżnia proces odwracalny
od nieodwracalnego. W samorzutnym przepływie ciepła
od ciała ciepłego do zimnego suma zmian entropii
jest dodatnia. Maszyna cieplna nie może zamienić
całego ciepła pobranego z jednego zbiornika na pracę
w cyklu. Przy obliczeniach entropii procesu
nieodwracalnego wybieramy dogodną hipotetyczną drogę
odwracalną między tymi samymi stanami.}
\topic{12}{Napięcie powierzchniowe, włoskowatość i ruchy Browna}{
Siły spójności między cząsteczkami powodują, że
utworzenie nowej powierzchni cieczy wymaga energii.
Napięcie powierzchniowe $\sigma$ jest energią na
jednostkę powierzchni albo siłą na jednostkę długości
w odpowiedniej geometrii. Dla kulistej kropli różnica
ciśnień wynosi $\Delta p=2\sigma/R$.

W cienkiej rurce kapilarnej wysokość podciągania
jest w przybliżeniu
$h=2\sigma\cos\theta/(\rho gr)$.
Znak $\cos\theta$ rozstrzyga, czy ciecz podnosi się,
czy opada. Ruchy Browna to losowe ruchy zawieszonych
drobnych cząstek pod wpływem uderzeń cząsteczek
płynu. Dla dyfuzji jednowymiarowej
$\langle x^2\rangle=2Dt$; jest to średnia
po wielu realizacjach, a nie tor pojedynczej cząstki.}
\newpage
\topic{13}{Transport w gazach: dyfuzja, ciepło, lepkość i prąd}{
Dyfuzja jest przepływem cząstek wywołanym gradientem
stężenia; pierwsze prawo Ficka zapisujemy
$\mathbf J=-D\nabla n$. Przewodnictwo cieplne
ma analogiczną postać $\mathbf q=-\kappa\nabla T$.
Lepkość przenosi pęd między warstwami gazu,
a dla prostego ścinania naprężenie jest proporcjonalne
do gradientu prędkości.

W kinetycznym obrazie współczynniki transportu zależą
od średniej drogi swobodnej i średniej szybkości cząstek.
Przewodnictwo elektryczne wymaga nośników ładunku:
$\mathbf J_e=\sigma_e\mathbf E$ w zakresie prawa Ohma.
Nie należy utożsamiać przewodnictwa elektrycznego
gazu z przewodnictwem cieplnym; pierwsze silnie zależy
od jonizacji. Relacje liniowe obowiązują blisko
równowagi i przy niezbyt dużych gradientach.}
\topic{14}{Potencjał, pole elektryczne, strumień i prawo Gaussa}{
Natężenie pola jest siłą na jednostkowy dodatni
ładunek próbny: $\mathbf E=\mathbf F/q$.
W elektrostatyce $\mathbf E=-\nabla\varphi$,
a różnica potencjałów odpowiada pracy pola
na jednostkę ładunku ze zmienionym znakiem.
Strumień przez powierzchnię to
$\Phi_E=\int_S\mathbf E\cdot d\mathbf S$.

Prawo Gaussa w próżni mówi
$\oint_S\mathbf E\cdot d\mathbf S
=Q_{\rm wew}/\varepsilon_0$.
Jest zawsze prawdziwe dla zamkniętej powierzchni,
ale praktycznie pozwala wyliczyć pole tylko przy
odpowiedniej symetrii, np. sferycznej, cylindrycznej
lub płaskiej. W postaci lokalnej
$\nabla\cdot\mathbf E=\rho/\varepsilon_0$.
Zewnętrzne ładunki mogą dawać pole na powierzchni,
lecz ich całkowity wkład do strumienia zamkniętego jest zerowy.}
\newpage
\topic{15}{Dielektryki, polaryzacja i zjawiska pokrewne}{
W dielektryku pole przesuwa ładunki związane lub
ustawia dipole. Wektor polaryzacji $\mathbf P$
jest momentem dipolowym na jednostkę objętości,
a $\mathbf D=\varepsilon_0\mathbf E+\mathbf P$.
Prawo Gaussa dla $\mathbf D$ obejmuje tylko ładunek
swobodny. W liniowym, izotropowym materiale
$\mathbf D=\varepsilon\mathbf E$.

Dla kondensatora płaskiego wypełnionego jednorodnym
dielektrykiem $C=\varepsilon A/d$, jeśli pomijamy
efekty brzegowe. Elektrostrykcja opisuje odkształcenie
wywołane polem; piezoelektryczność wiąże naprężenie
z polaryzacją w kryształach o odpowiedniej symetrii.
Ferroelektryki mają spontaniczną polaryzację
i pętlę histerezy. Trzeba odróżniać ich nieliniową
odpowiedź od prostego prawa $\mathbf P=\chi\varepsilon_0\mathbf E$.}
\topic{16}{Klasyczny model przewodnictwa elektrycznego}{
W modelu Drudego elektrony przewodnictwa są
nośnikami o średnim czasie relaksacji $\tau$.
W stałym polu ich średnia prędkość dryfu
spełnia $\mathbf v_d=-e\tau\mathbf E/m$.
Dla koncentracji $n$ prąd ma postać
$\mathbf J=-ne\mathbf v_d=\sigma\mathbf E$,
gdzie $\sigma=ne^2\tau/m$.

Model wyjaśnia prawo Ohma i jakościowo efekt Halla,
ale nie opisuje poprawnie wszystkich własności
metali, ponieważ pomija statystykę Fermiego
i strukturę pasmową. Czas relaksacji jest parametrem
efektywnym zależnym od rozpraszania na drganiach
sieci i defektach. Wartość dryfu jest mała
w porównaniu z typowymi mikroskopowymi szybkościami
elektronów; sygnał elektryczny rozchodzi się
znacznie szybciej niż pojedynczy nośnik.}
\newpage
\topic{17}{Prawo Ohma, SEM, Kirchhoff, praca i moc}{
Lokalne prawo Ohma to $\mathbf J=\sigma\mathbf E$.
Dla jednorodnego przewodnika o długości $L$
i przekroju $A$ mamy $R=\rho L/A$, a na jego
końcach $U=IR$. Siła elektromotoryczna źródła
jest pracą sił nieelektrostatycznych na jednostkę
ładunku w pełnym obiegu.

Pierwsze prawo Kirchhoffa wynika z zachowania
ładunku: suma prądów wpływających do węzła
równa się sumie wypływających. Drugie, w obwodzie
skupionym bez zmiennego strumienia magnetycznego,
sumuje napięcia w oczku do zera. Moc na oporniku
$P=UI=I^2R=U^2/R$, a energia wydzielona
w czasie $t$ przy stałej mocy wynosi $Pt$.
Przy zmiennym strumieniu trzeba uwzględnić
indukowaną SEM.}
\topic{18}{Prawo Ampère'a, siła przewodników i indukcja}{
W magnetostatyce prawo Ampère'a ma postać
$\oint\mathbf B\cdot d\boldsymbol\ell=\mu_0I_{\rm wew}$.
Pole długiego prostego przewodnika wynosi
$B=\mu_0I/(2\pi r)$. Dwa równoległe przewodniki
z prądami w tym samym kierunku przyciągają się;
siła na jednostkę długości to
$F/L=\mu_0I_1I_2/(2\pi d)$.

Prawo Faradaya mówi, że zmiana strumienia
magnetycznego indukuje siłę elektromotoryczną
$\mathcal E=-d\Phi_B/dt$. Znak minus wyraża
regułę Lenza: prąd indukowany przeciwdziała
zmianie strumienia. Strumień może się zmieniać
przez ruch obwodu, zmianę pola lub powierzchni.
Przy polu zależnym od czasu pełne prawo
Ampère'a wymaga dodatkowego członu Maxwella.}
\newpage
\topic{19}{Materia w polu magnetycznym}{
Magnetyzacja $\mathbf M$ to moment magnetyczny
na jednostkę objętości. Diamagnetyki mają
zwykle małą ujemną podatność: pole indukuje
moment przeciwny do pola. Paramagnetyki mają
dodatnią podatność, bo momenty atomowe
częściowo ustawiają się z polem.

Ferromagnetyki mogą mieć spontaniczne uporządkowanie
momentów poniżej temperatury Curie. Ich odpowiedź
jest nieliniowa i zależy od historii pola,
co daje histerezę oraz domeny magnetyczne.
W prostych materiałach liniowych
$\mathbf B=\mu_0(\mathbf H+\mathbf M)$
i $\mathbf M=\chi_m\mathbf H$, lecz nie należy
stosować stałej $\chi_m$ do całej pętli ferromagnetyka.
Mikroskopowe pochodzenie ferromagnetyzmu wymaga
kwantowej energii wymiany, której klasyczny
obraz małych magnesów nie wyjaśnia.}
\topic{20}{Drgania w obwodach LC i RLC}{
W idealnym obwodzie LC ładunek kondensatora
spełnia $L\ddot q+q/C=0$ i oscyluje z częstością
$\omega_0=1/\sqrt{LC}$. Energia przechodzi
między polem elektrycznym kondensatora
$q^2/(2C)$ a polem magnetycznym cewki
$Li^2/2$. W szeregowym RLC dochodzi tłumienie:
$L\ddot q+R\dot q+q/C=0$.

Przy sinusoidalnym wymuszeniu prądu zmiennego
moduł impedancji szeregowego obwodu wynosi
$|Z|=\sqrt{R^2+(\omega L-1/\omega C)^2}$.
Maksymalny prąd występuje przy
$\omega=1/\sqrt{LC}$ w tym idealizowanym modelu.
Dobroć $Q=\omega_0L/R$ opisuje ostrość rezonansu
przy małym tłumieniu. Trzeba podać, czy mówimy
o obwodzie szeregowym czy równoległym.}
\newpage
\topic{21}{Superpozycja fal, fale stojące i dudnienia}{
W ośrodku liniowym wychylenia fal sumują się.
Dwie fale o tej samej częstości i zgodnych fazach
mogą dać wzmocnienie, a o fazach przeciwnych
wygaszenie. Dwie przeciwbieżne fale o tej samej
amplitudzie tworzą falę stojącą,
$y=2A\sin(kx)\cos(\omega t)$.
Węzły pozostają nieruchome, a strzałki drgają
z największą amplitudą.

Przy bliskich częstościach $f_1$ i $f_2$
amplituda wypadkowa zmienia się wolno,
a częstość dudnień wynosi $|f_1-f_2|$.
W rurach i strunach warunki brzegowe wybierają
dozwolone długości fal i częstości własne.
Superpozycja zakłada liniowość równania falowego;
przy dużych amplitudach w ośrodku nieliniowym
proste dodawanie rozwiązań może zawieść.}
\topic{22}{Interferencja, doświadczenie Younga i Michelsona}{
Dla dwóch spójnych źródeł różnica dróg $\Delta$
określa różnicę faz $2\pi\Delta/\lambda$.
Wzmocnienie zachodzi przy $\Delta=m\lambda$,
a wygaszenie przy $\Delta=(m+1/2)\lambda$
dla źródeł zgodnych w fazie. W dalekim polu
dwu szczelin o odległości $d$ mamy
$\Delta\simeq d\sin\theta$.

W doświadczeniu Younga rozstaw prążków
na ekranie odległym o $L$ jest przy małych
kątach $\Delta x\simeq\lambda L/d$.
Natężenie dwóch jednakowych fal wynosi
$I=4I_0\cos^2(\Delta\phi/2)$.
Interferometr Michelsona rozdziela wiązkę
na dwa ramiona i łączy ją ponownie.
Przesunięcie zwierciadła o $\delta$ zmienia
drogę optyczną o $2\delta$, co umożliwia
precyzyjne pomiary długości i zmian refrakcji.}
\newpage
\topic{23}{Efekt Dopplera i fala uderzeniowa}{
Efekt Dopplera jest zmianą obserwowanej
częstotliwości wskutek ruchu źródła lub odbiorcy.
Dla dźwięku w nieruchomym ośrodku
$f'=f(v+v_o)/(v-v_s)$ przy odpowiednio
ustalonych dodatnich kierunkach zbliżania.
Wzór zależy od prędkości względem ośrodka.

Dla światła w próżni obowiązuje relatywistyczny
efekt Dopplera; przy oddalaniu wzdłuż linii
widzenia $f'=f\sqrt{(1-\beta)/(1+\beta)}$.
Nie wymaga on materialnego ośrodka.
Gdy źródło dźwięku przekracza prędkość fali,
fronty tworzą stożek Macha z
$\sin\theta=v_{\rm fala}/v_{\mathrm{zrodla}}$.
Przechodząca przez obserwatora fala uderzeniowa
powoduje skok ciśnienia. Nie jest to po prostu
przypadek graniczny zwykłego wzoru na ton
odbierany od poddźwiękowego źródła.}
\topic{24}{Widmo fal elektromagnetycznych}{
Fale elektromagnetyczne w próżni mają
$c=\lambda f$ i przenoszą energię; foton ma
$E=hf$. Od małych do dużych częstości
wyróżnia się fale radiowe, mikrofale,
podczerwień, światło widzialne, ultrafiolet,
promieniowanie rentgenowskie i gamma.
Granice pasm są umowne i częściowo zachodzą
na siebie.

Różnice praktyczne wynikają z energii fotonów
i oddziaływań z materią. Mikrofale silnie
oddziałują z pewnymi dipolami i strukturami,
widzialne pobudza receptory oka, a promieniowanie
UV, X oraz gamma może jonizować materię,
zależnie od energii. Rentgenowskie i gamma
często rozróżnia się także według sposobu
powstania: pierwsze wiąże się z elektronami,
drugie z przemianami jądrowymi.
Wszystkie te pasma podlegają równaniom Maxwella.}
\newpage

\topic{25}{Dyfrakcja światła i zasada Huygensa--Fresnela}{
Zasada Huygensa--Fresnela traktuje punkty czoła fali
jako źródła fal wtórnych, których superpozycja
wyznacza pole dalej. Dyfrakcja staje się widoczna,
gdy rozmiar przeszkody jest porównywalny
z długością fali. W obszarze Fraunhofera
obserwujemy rozkład w dalekim polu lub
w płaszczyźnie ogniskowej soczewki.

Dla szczeliny szerokości $a$ minima natężenia
spełniają $a\sin\theta=m\lambda$ dla
całkowitego $m\ne0$.
Centralne maksimum jest szersze niż boczne.
Dla siatki o stałej $d$ główne maksima
spełniają $d\sin\theta=m\lambda$.
Skończona liczba szczelin określa szerokość
prążków, a ich szerokość własna kształtuje
obwiednię. Równania zakładają spójne
oświetlenie i odpowiednią geometrię.}
\topic{26}{Optyka geometryczna i podstawowe przyrządy}{
Optyka geometryczna jest granicą falową,
w której długość fali jest mała względem
rozmiarów układu. Promień wskazuje kierunek
przenoszenia energii, a soczewka cienka
spełnia $1/f=1/p+1/q$ przy konsekwentnej
konwencji znaków. Powiększenie liniowe
wynosi $M=-q/p$.

Lupa zwiększa kąt widzenia małego przedmiotu;
jej przybliżone powiększenie kątowe zależy
od ogniskowej i punktu bliskiego oka.
Luneta łączy obiektyw tworzący obraz
dalekiego obiektu z okularem, który ten obraz
ogląda. Mikroskop najpierw tworzy powiększony
obraz pośredni obiektywem, a następnie
powiększa kątowo przez okular.
Rozdzielczości nie można zwiększać dowolnie:
ogranicza ją dyfrakcja i apertura układu.}
\newpage
\topic{27}{Odbicie, załamanie, całkowite odbicie i dyspersja}{
Na granicy ośrodków kąt odbicia równa się
kątowi padania. Prawo Snelliusa
$n_1\sin\theta_1=n_2\sin\theta_2$
wynika z ciągłości fazy wzdłuż granicy.
Przy przejściu z większego $n$ do mniejszego
całkowite wewnętrzne odbicie zachodzi
dla $\theta_1>\theta_c$, gdzie
$\sin\theta_c=n_2/n_1$.

Współczynnik załamania zależy od częstotliwości.
Dyspersja normalna oznacza zwykle wzrost $n$
z częstością w obszarze dalekim od rezonansu;
w pobliżu pasma absorpcyjnego może wystąpić
dyspersja anomalna. Szybkość fazowa $c/n$
nie musi równać się grupowej. W światłowodzie
całkowite odbicie pomaga prowadzić światło,
ale rzeczywiste straty zależą też od absorpcji
i rozpraszania.}
\topic{28}{Polaryzacja, prawo Malusa i dwójłomność}{
Polaryzacja opisuje kierunek drgań wektora
pola elektrycznego poprzecznej fali.
Idealny polaryzator przepuszcza składową
wzdłuż swojej osi; dla światła liniowo
spolaryzowanego prawo Malusa daje
$I=I_0\cos^2\theta$. Dla światła
niespolaryzowanego pierwszy idealny
polaryzator przepuszcza średnio połowę
natężenia.

Przy padaniu pod kątem Brewstera światło
odbite może być liniowo spolaryzowane;
$\tan\theta_B=n_2/n_1$ dla ośrodków
niemagnetycznych. Dwójłomny kryształ
ma różne współczynniki załamania dla
różnych polaryzacji. Dwójłomność może
też być wymuszona naprężeniem lub polem.
Aktywność optyczna obraca płaszczyznę
polaryzacji i wiąże się m.in. z chiralnością
ośrodka.}
\newpage
\topic{29}{Naturalna szerokość linii i stany metatrwałe}{
Stan wzbudzony o skończonym czasie życia
$\tau$ ma niezerową szerokość energetyczną.
Dla zaniku wykładniczego profil linii
jest lorentzowski, a rząd wielkości
szerokości energii to $\Delta E\sim\hbar/\tau$;
precyzyjny współczynnik zależy od przyjętej
definicji szerokości. Krótki czas życia
oznacza szerszą linię.

W doświadczeniu obserwowana szerokość może
być większa z powodu efektu Dopplera,
zderzeń lub aparatury. Stan metatrwały
ma wyjątkowo długi czas życia, ponieważ
przejście do niższego stanu jest słabo
dozwolone przez reguły wyboru lub
wymaga wyższego rzędu procesu.
Takie stany umożliwiają gromadzenie obsadzenia
w laserach. Trzeba odróżnić naturalną
szerokość pojedynczego przejścia od
poszerzenia całego badanego widma.}
\topic{30}{Doświadczenie Rutherforda i model atomu}{
Rutherford badał rozpraszanie cząstek alfa
na cienkiej folii. Większość przelatywała
z niewielkim odchyleniem, ale mała część
ulegała dużym kątom rozproszenia.
Model dodatniego ładunku rozłożonego
równomiernie w całej objętości atomu
nie wyjaśniał łatwo tych rzadkich
silnych odchyleń.

Wniosek był taki, że dodatni ładunek
i niemal cała masa są skupione w bardzo
małym jądrze. Klasyczny wzór Rutherforda
na różniczkowy przekrój czynny ma zależność
$d\sigma/d\Omega\propto\sin^{-4}(\theta/2)$
dla odpychania kulombowskiego cząstek
punktowych. Model jądrowy rozwiązywał
problem rozpraszania, lecz klasyczna
elektrodynamika nie wyjaśniała stabilności
orbit elektronowych; wymagała ona
opisu kwantowego.}
\newpage
\topic{31}{Widmo promieniowania rentgenowskiego}{
Promieniowanie rentgenowskie powstaje
gdy szybkie elektrony hamują w materiale
anody albo wybijają elektrony z wewnętrznych
powłok atomu. Hamowanie daje widmo ciągłe.
Jego graniczna długość fali spełnia
$\lambda_{\min}=hc/(eU)$, gdy elektron
przyspieszono napięciem $U$ i całą
energię oddał jednemu fotonowi.

Promieniowanie charakterystyczne tworzy
linie związane z przejściami między
powłokami po powstaniu wakatu.
Ich energie zależą od pierwiastka anody,
podczas gdy granica widma ciągłego
zależy od napięcia. Intensywność
i kształt widma zależą również od
filtracji i absorpcji w materiale.
Przy opisie widma trzeba rozróżnić
proces wytwarzania fotonów od ich
późniejszego osłabienia w ośrodku.}
\topic{32}{Wzbudzanie i deekscytacja jąder oraz atomów}{
Wzbudzone jądro może oddać energię
jako kwant gamma albo przez konwersję
wewnętrzną: energia przejścia trafia
bezpośrednio do elektronu atomowego.
Po powstaniu wakatu w powłoce atom
może wyemitować promieniowanie X
lub elektron Augera. Wybór kanału
zależy od energii i struktury atomu.

Efekt Mössbauera polega na niemal
bezodrzutowej emisji i absorpcji
kwantów gamma przez jądra osadzone
w sieci krystalicznej. Cała sieć
przejmuje pęd, więc energia odrzutu
pojedynczego jądra staje się znikoma.
Umożliwia to bardzo wąskie rezonanse
i badanie subtelnych przesunięć
energetycznych. Należy odróżniać
konwersję wewnętrzną od rozpadu beta:
w pierwszym przypadku jądro
nie zmienia liczby protonów.}
\newpage
\topic{33}{Jądro atomowe: odkrycie, budowa i własności}{
Doświadczenie Rutherforda wskazało
istnienie małego, dodatnio naładowanego
jądra. Składa się ono z protonów
i neutronów, czyli nukleonów.
Liczba protonów $Z$ określa pierwiastek,
a liczba masowa $A=Z+N$ sumę nukleonów.
Izotopy mają to samo $Z$, lecz różne $N$.

Typowy promień jądra przybliża wzór
$R=r_0A^{1/3}$, gdzie $r_0$ jest rzędu
femtometra. Oznacza to prawie stałą gęstość
materii jądrowej. Masa związanego jądra
jest mniejsza od sumy mas swobodnych
nukleonów; defekt masy odpowiada
energii wiązania $E_b=\Delta mc^2$.
Jądra mają spin, moment magnetyczny
i dyskretne stany wzbudzone.
Stabilność wynika z równowagi sił
jądrowych i odpychania Coulomba.}
\topic{34}{Model kroplowy i powłokowy jądra}{
Model kroplowy traktuje jądro
jak naładowaną kroplę o niemal stałej
gęstości. Półempiryczny wzór masowy
zawiera składniki objętościowy,
powierzchniowy, Coulomba, asymetrii
i parowania. Pomaga wyjaśnić
ogólne trendy energii wiązania
oraz możliwość rozszczepienia
ciężkich jąder.

Model powłokowy opisuje nukleony
poruszające się w średnim potencjale
z uwzględnieniem zakazu Pauliego.
Silne sprzężenie spin--orbita
prowadzi do liczb magicznych
i zwiększonej stabilności
zamkniętych powłok. Model
wyjaśnia spiny i momenty wielu
jąder, których sama kropla
nie opisuje. Oba modele
są przybliżeniami o różnych
obszarach skuteczności;
jądro może wykazywać zarazem
cechy kolektywne i indywidualne.}
\newpage
\topic{35}{Spontaniczne przemiany jądrowe}{
Rozpad alfa emituje jądro helu i jest
istotny dla ciężkich nuklidów;
bariera Coulomba jest pokonywana
przez tunelowanie. W rozpadzie beta
oddziaływanie słabe zmienia neutron
w proton lub odwrotnie, z emisją
leptonów zapewniającą zachowanie
ładunku i liczby leptonowej.
Rozpad gamma usuwa energię wzbudzenia
bez zmiany $A$ i $Z$.

Prawo zaniku pojedynczego kanału
ma postać $N(t)=N_0e^{-\lambda t}$,
a czas połowicznego zaniku
$T_{1/2}=\ln2/\lambda$.
Różne kanały mają częściowe
stałe zaniku sumujące się
do całkowitej. Energia wydzielona
$Q$ wynika z różnicy mas
początkowych i końcowych.
Przemiana jest spontaniczna,
gdy bilans energii pozwala
na stan końcowy, ale szybkość
może być bardzo mała.}
\topic{36}{Oddziaływanie promieniowania jądrowego z materią}{
Cząstki naładowane tracą energię
głównie przez jonizację i wzbudzanie
atomów. Ciężkie cząstki alfa
mają krótki zasięg i gęsto
jonizują; elektrony beta ulegają
też silnemu rozpraszaniu
i mogą wytwarzać promieniowanie
hamowania. Fotony gamma oddziałują
m.in. przez efekt fotoelektryczny,
rozpraszanie Comptona i tworzenie
par przy odpowiedniej energii.

Neutrony są elektrycznie obojętne,
więc oddziałują z jądrami
przez rozpraszanie i reakcje
jądrowe. Osłabienie wąskiej
wiązki fotonów bywa opisane
$I(x)=I_0e^{-\mu x}$,
ale współczynnik zależy od
energii i materiału. Zasięg
cząstki i współczynnik osłabienia
fotonów to różne wielkości;
przy ochronie radiologicznej
trzeba uwzględnić również
promieniowanie wtórne.}
\newpage
\topic{37}{Rozszczepienie jąder i energetyka jądrowa}{
Ciężkie jądro może rozpaść się
na dwa lżejsze fragmenty.
W rozszczepieniu indukowanym
neutronem powstają zwykle
dodatkowe neutrony, co pozwala
na reakcję łańcuchową.
Energia wydzielona wynika
z większej energii wiązania
na nukleon w produktach:
$Q=(m_{\rm przed}-m_{\rm po})c^2$.

Reaktor utrzymuje średnio
jedną następną reakcję
z każdego rozszczepienia
w stanie krytycznym.
Moderator spowalnia neutrony
w reaktorach termicznych,
a pręty sterujące zmieniają
ich pochłanianie. Ciepło
z reakcji jest przenoszone
do obiegu energetycznego
i zamieniane na pracę
przez turbinę. Bilans
techniczny obejmuje sprawność
cieplną, chłodzenie,
odpady promieniotwórcze
i bezpieczeństwo układu.}
\newpage

\section*{II. Fizyka teoretyczna}
\topic{1}{Lagrangian i oscylator harmoniczny}{
Lagrangian to różnica energii kinetycznej i potencjalnej
$L(q,\dot q,t)=T-U$. Zasada stacjonarnego działania
$S=\int L\,dt$ prowadzi do równania Eulera--Lagrange'a
$d(\partial L/\partial\dot q)/dt-\partial L/\partial q=0$.
Dla masy na sprężynie
$L=m\dot x^2/2-kx^2/2$, zatem
$m\ddot x+kx=0$.

Rozwiązanie ogólne to
$x(t)=A\cos(\omega_0t)+B\sin(\omega_0t)$,
gdzie $\omega_0=\sqrt{k/m}$.
Warunki $x(0)=x_0$ i $\dot x(0)=v_0$
dają $A=x_0$, $B=v_0/\omega_0$.
Wprowadzenie współrzędnych uogólnionych
pozwala uwzględnić więzy bez jawnego
rozpisywania sił reakcji. Ważne jest,
aby w równaniu ruchu uwzględnić siły
niezachowawcze osobno jako siły uogólnione,
jeśli nie można ich zawrzeć w potencjale.}
\topic{2}{Stopnie swobody i ruch na sferze}{
Swobodny punkt w przestrzeni ma trzy stopnie
swobody. Więz $x^2+y^2+z^2=R^2$ usuwa
jeden, pozostawiając dwa: kąty $\theta,\phi$.
Dla punktu o masie $m$ poruszającego się
bez tarcia po nieruchomej sferze
$L=mR^2(\dot\theta^2+\sin^2\theta\dot\phi^2)/2$.

Równanie dla $\phi$ daje
$d(mR^2\sin^2\theta\dot\phi)/dt=0$,
ponieważ $\phi$ jest współrzędną cykliczną.
Dla $\theta$ otrzymujemy
$\ddot\theta-\sin\theta\cos\theta\,\dot\phi^2=0$.
Zachowany pęd sprzężony z $\phi$
jest składową momentu pędu względem
osi $z$. W tym przykładzie siła
reakcji więzu nie wykonuje pracy
w dozwolonych przesunięciach
wirtualnych, dlatego nie pojawia
się w równaniach dla kątów.}
\newpage
\topic{3}{Hamiltonian i oscylator harmoniczny}{
Pęd kanoniczny definiujemy jako
$p=\partial L/\partial\dot q$, a Hamiltonian
$H=p\dot q-L$ po wyeliminowaniu prędkości.
Dla oscylatora $p=m\dot x$ oraz
$H=p^2/(2m)+kx^2/2$.
Równania Hamiltona to
$\dot x=\partial H/\partial p=p/m$
i $\dot p=-\partial H/\partial x=-kx$.

Ich połączenie daje
$\ddot x+\omega_0^2x=0$.
W przestrzeni fazowej trajektorie
o stałej energii są elipsami:
$p^2/(2m)+kx^2/2=E$.
Dla Hamiltonianu niezależnego
jawnie od czasu energia jest
stała. Formalizm Hamiltona
opisuje ewolucję punktu
$(x,p)$, co ułatwia analizę
symetrii, zachowanych wielkości
i przejście do mechaniki
statystycznej oraz kwantowej.}
\topic{4}{Wahadło matematyczne w opisie Hamiltona}{
Jedyną współrzędną wahadła
o długości $\ell$ jest kąt $\theta$.
Dla masy punktowej $m$ i potencjału
liczonego od dolnego położenia
$L=m\ell^2\dot\theta^2/2
-mg\ell(1-\cos\theta)$.
Pęd kanoniczny wynosi
$p_\theta=m\ell^2\dot\theta$.

Hamiltonian jest sumą
$H=p_\theta^2/(2m\ell^2)
+mg\ell(1-\cos\theta)$.
Równania Hamiltona:
$\dot\theta=p_\theta/(m\ell^2)$
i $\dot p_\theta=-mg\ell\sin\theta$.
W przybliżeniu $|\theta|\ll1$ mamy
$\sin\theta\simeq\theta$,
więc $\ddot\theta+(g/\ell)\theta=0$
i okres $T\simeq2\pi\sqrt{\ell/g}$.
Przy dużych amplitudach okres
zależy od amplitudy; przybliżenie
małych drgań przestaje działać.}
\newpage
\topic{5}{Ładunek próbny i siła Lorentza}{
Ładunek próbny jest na tyle mały,
że nie zmienia istotnie badanego pola.
Na ładunek $q$ poruszający się
z prędkością $\mathbf v$ działa
siła Lorentza
$\mathbf F=q(\mathbf E+\mathbf v\times\mathbf B)$.
Stałe jednorodne pole elektryczne
przyspiesza cząstkę w kierunku
$q\mathbf E$; bez pola magnetycznego
tor jest paraboliczny przy
odpowiednich warunkach początkowych.

Samo pole magnetyczne nie wykonuje
pracy, bo $(\mathbf v\times\mathbf B)
\cdot\mathbf v=0$.
Składowa prędkości prostopadła
do $\mathbf B$ daje ruch kołowy
o promieniu $r=m v_\perp/(|q|B)$
w granicy nierelatywistycznej,
a równoległa ruch jednostajny.
Złożenie tych ruchów tworzy
helisę. Znak ładunku odwraca
kierunek zakrzywienia toru.}
\topic{6}{Prawo Gaussa dla pól E i B}{
Prawo Gaussa dla elektryczności
w postaci całkowej to
$\oint_S\mathbf E\cdot d\mathbf S
=Q_{\rm wew}/\varepsilon_0$,
a lokalnie
$\nabla\cdot\mathbf E=\rho/\varepsilon_0$.
Dla pola magnetycznego
$\oint_S\mathbf B\cdot d\mathbf S=0$
i $\nabla\cdot\mathbf B=0$;
strumień wchodzący i wychodzący
z zamkniętej powierzchni
bilansuje się.

Dla punktowego ładunku w próżni
symetria sferyczna daje
$E(4\pi r^2)=q/\varepsilon_0$,
stąd prawo Coulomba
$\mathbf E=q\hat{\mathbf r}/
(4\pi\varepsilon_0r^2)$.
Zastosowanie wymaga, by na
wybranej sferze wartość $E$
była stała i pole radialne.
Prawo magnetyczne opisuje
brak obserwowanych izolowanych
monopoli w klasycznych
równaniach Maxwella.}
\newpage
\topic{7}{Prawo Faradaya i indukowana SEM}{
Prawo Faradaya w postaci całkowej
mówi $\oint_C\mathbf E\cdot
d\boldsymbol\ell=-d\Phi_B/dt$
dla nieruchomego konturu $C$,
gdzie $\Phi_B=\int_S\mathbf B
\cdot d\mathbf S$.
Lokalnie jest to
$\nabla\times\mathbf E
=-\partial\mathbf B/\partial t$.
Zmienne pole magnetyczne
wytwarza wirowe pole elektryczne.

Dla poruszającego się obwodu
SEM zawiera także człon
z siłą magnetyczną
$\mathbf v\times\mathbf B$.
W generatorze obrót ramki
zmienia strumień i wywołuje
okresową SEM. Znak w prawie
Faradaya wyraża regułę Lenza:
prąd indukowany przeciwdziała
zmianie strumienia, która go
wywołała. Rozróżnienie
nieruchomej i ruchomej
granicy całkowania zapobiega
błędnej interpretacji
samego członu $\partial B/\partial t$.}
\topic{8}{Prawo Ampère'a--Maxwella i kondensator}{
Pełne prawo ma postać
$\oint_C\mathbf B\cdot
d\boldsymbol\ell
=\mu_0(I_{\rm przew}
+\varepsilon_0\,d\Phi_E/dt)$.
W postaci lokalnej:
$\nabla\times\mathbf B
=\mu_0\mathbf J
+\mu_0\varepsilon_0
\partial\mathbf E/\partial t$.
Człon z pochodną pola
elektrycznego nazywa się
prądem przesunięcia.

Podczas ładowania kondensatora
prądem stałym przewodzenia $I$
przez przewód, między okładkami
nie płynie taki sam prąd
nośników przez próżnię.
Pole $E$ jednak narasta
i dla idealnego kondensatora
$\varepsilon_0\,d\Phi_E/dt=I$.
Dzięki temu wynik prawa
Ampère'a nie zależy od
wyboru powierzchni rozpiętej
na tym samym konturze.
Poprawka jest zgodna
z lokalnym zachowaniem ładunku.}
\newpage

\input{teoretyczna_9_17}
\section*{III. Ekonofizyka}
\topic{1}{Konstrukcja i analiza portfela papierów wartościowych}{
Portfel opisują wagi $w_i$ poszczególnych aktywów,
zwykle spełniające $\sum_iw_i=1$.
Jego oczekiwana stopa zwrotu to
$\mu_p=\mathbf w^\top\boldsymbol\mu$,
a wariancja
$\sigma_p^2=\mathbf w^\top\Sigma\mathbf w$,
gdzie $\Sigma$ jest macierzą kowariancji
zwrotów. Korelacje decydują,
ile ryzyka można ograniczyć
przez dywersyfikację.

W ujęciu Markowitza wybiera się
portfele efektywne: dla danego
ryzyka mają najwyższy oczekiwany
zwrot albo dla danego zwrotu
najmniejsze ryzyko.
Analiza wymaga oszacowania
$\boldsymbol\mu$ i $\Sigma$
z danych, określenia horyzontu,
kosztów transakcyjnych
oraz ograniczeń wag.
Wynik może być niestabilny,
ponieważ błędy oszacowania
średnich zwrotów są duże;
warto sprawdzać odporność
rozwiązania na zmianę
okna danych.}
\topic{2}{Portfel najmniejszego ryzyka i portfel optymalny}{
Portfel minimum wariancji
rozwiązuje problem
$\min_{\mathbf w}\mathbf w^\top
\Sigma\mathbf w$ przy
$\mathbf 1^\top\mathbf w=1$.
Bez zakazu krótkiej sprzedaży
i przy odwracalnej $\Sigma$
rozwiązanie ma postać
$\mathbf w_{\rm GMV}
=\Sigma^{-1}\mathbf 1/
(\mathbf 1^\top\Sigma^{-1}\mathbf 1)$.

Portfel optymalny wymaga
dodatkowego kryterium
preferencji inwestora,
na przykład maksymalizacji
$\mathbf w^\top\boldsymbol\mu
-\lambda\mathbf w^\top\Sigma\mathbf w/2$.
Parametr $\lambda$ wyraża
awersję do ryzyka.
Portfel minimum wariancji
nie musi maksymalizować
tej funkcji. Dodatkowe
ograniczenia, koszty
i niedokładność estymacji
mogą znacznie zmienić
wagi, dlatego rozwiązanie
algebraiczne jest punktem
wyjścia do analizy,
a nie prognozą wyników.}
\newpage
\topic{3}{Korelacje, zwroty i taksonomia portfela}{
Z cen $P_i(t)$ oblicza się
zwroty proste lub logarytmiczne,
np. $r_i(t)=\ln P_i(t)-
\ln P_i(t-\Delta t)$.
Korelacja Pearsona
$\rho_{ij}$ mierzy liniową
współzmienność zwrotów,
nie samych poziomów cen.
Korelacje cen często bywają
pozorne z powodu trendów.

Odległość
$d_{ij}=\sqrt{2(1-\rho_{ij})}$
pozwala zbudować drzewo
minimalne lub hierarchię
podobieństw między aktywami.
Taka taksonomia może ujawnić
grupy sektorowe i wspólne
czynniki, ale zależy
od okna czasowego,
częstotliwości danych
i sposobu estymacji.
Korelacja nie mierzy
zależności przyczynowej
ani dobrze nie opisuje
wszystkich ekstremalnych
współruchów. Przed
porównaniem serii trzeba
zsynchronizować notowania
i obsłużyć brakujące dane.}
\topic{4}{Krachy jako analogia przejść fazowych}{
W ekonofizyce krach bywa
modelowany jako gwałtowna
zmiana stanu układu
złożonego z oddziałujących
uczestników. Analogią
przejścia fazowego są
narastające korelacje,
niestabilność i możliwe
zachowanie skalowe.
Modele logarytmiczno-
periodyczne próbują
opisać przyspieszanie
bańki przed punktem
krytycznym.

Nie ma jednak uniwersalnego
prawa fizycznego,
które wyznacza datę krachu.
Dopasowanie modelu
do historii może łatwo
przeuczyć dane,
a wybór początku bańki
zmienia wynik.
Rzetelne badanie wymaga
definicji krachu,
testów poza próbą,
przedziałów niepewności
i porównania z prostszymi
modelami. Wskaźnik
ostrzegawczy jest
hipotezą statystyczną,
nie pewną zapowiedzią
przyszłego zdarzenia.}
\newpage
\topic{5}{Statystyczne własności rozkładów zwrotów}{
Zwroty finansowe mają
często rozkład o grubych
ogonach względem rozkładu
normalnego. Duże odchylenia
występują częściej niż
przewiduje model Gaussa.
Znaki krótkookresowych
zwrotów mogą mieć słabą
autokorelację, podczas
gdy wartości bezwzględne
lub kwadraty zwrotów
wykazują klastrowanie
zmienności.

Analiza obejmuje średnią,
wariancję, skośność,
kurtozę i funkcję
autokorelacji, ale
estymatory zależą
od okresu i rozdzielczości
danych. Rozkłady zwrotów
dodatkowo zmieniają się
w czasie, więc jedna
funkcja gęstości nie musi
opisywać całego rynku.
Trzeba sprawdzić wpływ
ekstremów, świąt,
braku płynności
i zmian metodologii.
Ciężki ogon sam w sobie
nie dowodzi konkretnego
mechanizmu generowania
cen.}
\topic{6}{Modele AR, MA i ARMA a rozkład Gaussa}{
Model AR($p$) opisuje
bieżący zwrot przez
poprzednie zwroty
i innowację:
$r_t=c+\sum_{i=1}^p
\phi_i r_{t-i}+\varepsilon_t$.
Model MA($q$) używa
bieżącej i poprzednich
innowacji, a ARMA($p,q$)
łączy oba składniki.
Parametry opisują
strukturę warunkowej
średniej szeregu.

Jeśli innowacje są
niezależne i normalne,
stacjonarny liniowy
model ARMA ma
gaussowskie rozkłady
skończenie wymiarowe.
Sam zapis ARMA
nie wymusza jednak
normalności; można
użyć innowacji
o grubych ogonach.
Model dla średniej
nie wystarcza
do opisu zmiennej
w czasie wariancji.
Dobór rzędów
opiera się na
diagnostyce reszt,
informacji kryterialnej
i sprawdzeniu
stacjonarności.}
\newpage
\topic{7}{Modele ARCH i GARCH}{
ARCH zakłada, że
warunkowa wariancja
zależy od przeszłych
kwadratów innowacji,
np. $h_t=\omega+
\alpha\varepsilon_{t-1}^2$.
GARCH dodaje poprzednią
wariancję:
$h_t=\omega+
\alpha\varepsilon_{t-1}^2+
\beta h_{t-1}$.
Zapis
$\varepsilon_t=\sqrt{h_t}z_t$
oddziela zmienną
skalę od standaryzowanej
innowacji $z_t$.

Dla prostego GARCH(1,1)
przy $\omega>0$,
$\alpha,\beta\geq0$
i $\alpha+\beta<1$
istnieje skończona
wariancja bezwarunkowa
$\omega/(1-\alpha-\beta)$.
Model uchwytuje
klastrowanie zmienności,
ale klasyczna wersja
symetrycznie reaguje
na dodatnie i ujemne
szoki tej samej wielkości.
W praktyce sprawdza się
reszty standaryzowane,
stabilność parametrów
i jakość prognoz
poza próbą.}
\topic{8}{Skalowanie, fraktalność i multifraktalność}{
Skalowanie oznacza
systematyczną zależność
statystyk przy zmianie
skali czasu. Dla
prostego procesu
samopodobnego moment
rzędu $q$ może spełniać
$E|r_{\Delta t}|^q
\propto(\Delta t)^{\zeta(q)}$.
W monofraktalnym
przybliżeniu $\zeta(q)$
jest liniowe w $q$;
multifraktalność
wiąże się z
nieliniowością tej
funkcji.

Szeregi finansowe
mogą pokazywać
zależność zmienności
na wielu skalach.
Modele multifraktalne
próbują połączyć
grube ogony,
klastrowanie
i skalowanie.
Estymacja jest jednak
wrażliwa na długość
próby, zakres skal
i skrajne obserwacje.
Pozorna multifraktalność
może też wynikać
ze zmian reżimu.
Wnioski wymagają
porównania z
danymi losowo
przemieszanymi
i modelami
odniesienia.}
\newpage
\section*{Podstawowe źródła do części III}
\begin{itemize}
\item H. Markowitz,
\href{https://onlinelibrary.wiley.com/doi/10.1111/j.1540-6261.1952.tb01525.x}
{Portfolio Selection}, 1952.
\item R. N. Mantegna,
\href{https://arxiv.org/abs/cond-mat/9802256}
{Hierarchical Structure in Financial Markets}, 1998.
\item R. F. Engle,
\href{https://web.pdx.edu/~crkl/ec572/readings/Engle82.pdf}
{Autoregressive Conditional Heteroscedasticity}, 1982.
\item T. Bollerslev,
\href{https://www.sciencedirect.com/science/article/pii/0304407686900631}
{Generalized Autoregressive Conditional Heteroskedasticity}, 1986.
\item E. Bacry, J. Delour, J.-F. Muzy,
\href{https://arxiv.org/abs/cond-mat/0005405}
{A multifractal random walk}, 2000.
\end{itemize}

\end{document}
